\section{Open questions}

The following five questions were left unresolved by this replication.
Each is scoped so that a concrete follow-up experiment could close it
without leaving the free-endpoint, small-instance regime of the QC-100
wave.

\begin{enumerate}
\item \textbf{Non-Pauli noise robustness.}
Does the efficient-ZNE (2-scale linear + random-gate folding) recipe
retain its advantage over full ZNE (3-scale Richardson + global folding)
under coherent over-rotations on CNOT and biased amplitude/phase damping?
\emph{Basis:} we tested only two-qubit depolarising + amplitude damping
(the paper's Fig.~2 recipe); the Pauli-Lindblad expansion that underlies
Richardson-style ZNE is known to degrade under coherent noise, and the
paper does not enumerate results in that regime.
\emph{Next steps:} extend \texttt{zne\_reproduction.py} with a coherent
$R_{zz}(\delta)$ channel after each \texttt{cx} for
$\delta \in \{0.005, 0.01, 0.02\}$~rad and a biased $T_2/T_1$ ratio in
$\{0.5, 1.0, 2.0\}$; re-run the CNOT sweep and 30-trial precision study;
declare the efficient variant robust if MAE(eff)/MAE(full) stays
$\le 1.1$ across all settings.

\item \textbf{Composition with virtual distillation.}
Does composing efficient ZNE with virtual distillation (Huggins et al.,
2021) give a strictly better bias-variance product than either alone at
fixed shot budget?
\emph{Basis:} the two mechanisms attack different error components (bias
in $\varepsilon$ vs off-diagonal state components), but their composition
on the paper's Fig.~2 target has not been demonstrated, and VD amplifies
shot noise as $\sim 1/M^2$.
\emph{Next steps:} implement $M=2$ virtual distillation on
$\ket{11}\bra{11}$ with two independent copies (4 qubits total); apply
Mitiq's efficient-ZNE on top; compare
$(\text{bias}, \text{variance}, \text{shots})$ for
$\{\text{raw}, \text{VD}, \text{eff-ZNE}, \text{VD+eff-ZNE}\}$ at fixed
$24\,576$ shots.

\item \textbf{ML-selected extrapolation schedule.}
Can a learned selector (choosing scale factors and extrapolator from the
observed raw survival $\Pr(\ket{11})$) beat both fixed Richardson-3 and
fixed Linear-2 across a distribution of noise magnitudes and depths?
\emph{Basis:} both arms here use fixed scale-factor sets; a data-driven
selector could keep full-ZNE accuracy where needed and drop to a cheap
2-scale linear when the raw signal is already high, saving shots.
\emph{Next steps:} generate $500$ $(\varepsilon, T_1, n_c)$ triples with
per-triple ground truth and MAE for six candidate schedules; train a
gradient-boosted classifier on
$(\text{raw survival}, \sqrt{\text{shots}})$; evaluate held-out MAE
against the fixed baselines; target $\ge 10\%$ MAE reduction at equal
shots.

\item \textbf{Noise-model-adaptive schedule from device calibration.}
How does a schedule chosen from a device-calibration report compare to
the fixed Linear-2 recipe on Fig.~2/3 under realistic IBM Falcon-era
calibration data (per-pair $\varepsilon_{cx}$, per-qubit $T_1$, $T_2$)?
\emph{Basis:} the paper's fixed scale factors are optimal at a single
$(\varepsilon, T_1)$ point; an adaptive schedule that reads calibration
should Pareto-dominate any fixed schedule when calibration varies across
the device.
\emph{Next steps:} pull a public IBM~Q calibration snapshot for $\sim 10$
qubit pairs; simulate Fig.~2/3 under each pair's noise; compare fixed
Linear-2 against $s = (1, 1 + k\varepsilon_{cx})$ for
$k \in \{50, 100, 200\}$; target $\ge 15\%$ median-MAE improvement.

\item \textbf{Beyond the toy circuit: RB, mirror, brick-layer benchmarks.}
Does efficient ZNE preserve its advantage on canonical hardware-benchmark
circuits (randomized benchmarking, mirror benchmarking, layered QV), and
how does the crossover CNOT count (where efficient no longer beats full)
scale with circuit width?
\emph{Basis:} the paper's headline example is a 2-qubit single-observable
circuit; practitioners care about many-qubit benchmarks where per-CNOT
noise composes non-trivially and where the ``observable near $1$''
assumption is replaced by a decaying survival probability.
\emph{Next steps:} add three benchmark families ---
(a) 4-qubit RB at lengths $\{10, 20, 40\}$,
(b) 4-qubit mirror circuits (Proctor et al.),
(c) 6-qubit brick-layer at depths $\{5, 10, 20\}$ ---
run raw / full-ZNE / eff-ZNE with the current noise recipe, record MAE
and shots-to-target-precision, and identify the $n_c$ crossover as a
function of qubit count.
\end{enumerate}
